\begin{textsample}{POS Dim 4 – pa – Score 13.00 – pa/pa\_ef\_29.txt}  \label{ex:f4_pos_011}
EDUCAÇÃO DAS RELAÇÕES ÉTNICO-RACIAIS E QUILOMBOLAS A política curricular proposta para a educação das relações étnico-raciais e quilombola deve estar fundada em dimensões históricas, sociais, \textbf{antropológicas} oriundas da realidade brasileira, buscando combater o \textbf{racismo} e \textbf{discriminações} que atingem negros e \textbf{índios}.
A referida proposta apresenta como metas o direito dos negros e dos \textbf{índios} de se reconhecerem na cultura nacional, de expressarem visões próprias de mundo, de se manifestarem com autonomia individual e coletivamente.
Esse direito garante o acesso dos referidos cidadãos a cursarem todos os níveis de ensino em escolas devidamente instaladas e equipadas, orientados por professores com formação para lidar com as relações produzidas pelo \textbf{racismo} e \textbf{discriminações}, capazes de conduzir a reeducação das relações entre diferentes grupos étnico-raciais e a valorização da história, da cultura e da \textbf{identidade} dos \textbf{indígenas} e descendentes de africanos.
Nela se propõe a divulgação e a produção de conhecimentos, a formação de atitudes, posturas e valores que eduquem cidadãos orgulhosos de seu pertencimento étnico-racial para interagirem na construção de uma nação democrática, justa e inclusiva em que todos igualmente tenham seus direitos garantidos e sua \textbf{identidade} valorizada.
A política curricular para a educação das relações étnico-raciais e quilombola tem como princípios : Esse princípio deve conduzir : - À igualdade básica da pessoa humana como sujeito de direitos ; - À compreensão de que a sociedade é formada por pessoas que pertencem a grupos étnico-raciais distintos, que possuem cultura e história próprias, igualmente valiosas e que em conjunto constroem, na nação brasileira, sua história ; - Ao conhecimento e à valorização da história dos povos africanos e da cultura afro-brasileira na construção histórica e cultural brasileira ; - À superação da indiferença, injustiça e desqualificação com que os negros, os povos \textbf{indígenas} e também as classes \textbf{populares} às quais os negros, no geral, pertencem, são comumente tratados ; - À desconstrução ideologia do branqueamento por meio de questionamentos e análises críticas, objetivando eliminar conceitos, ideias, comportamentos veiculados pelo mito da democracia \textbf{racial} que tanto mal faz a negros, a \textbf{índios} e a \textbf{brancos} ; - À análise das relações étnico-raciais e sociais com o estudo de história e cultura afro-brasileira, africana e \textbf{indígena}, de informações e subsídios que permitam aos profissionais da educação formular concepções e percursos pedagógicos pautados na superação de preconceitos e com capacidade de construir posturas e atitudes respeitosas ; - Ao estabelecimento de uma relação dialógica entre os diferentes, com a finalidade de negociações visando à construção de uma sociedade justa. 5.6.2 Fortalecimento de \textbf{Identidades} e de Direitos Esse princípio deve orientar para : - O desencadeamento de processos de afirmação de \textbf{identidades}, de historicidade \textbf{negada} ou distorcida ; - O rompimento com imagens negativas forjadas por diferentes meios de comunicação, contra os negros e os povos \textbf{indígenas} ; - O respeito à diversidade identitária em \textbf{contraposição} à ideia de uma \textbf{identidade} humana universal ; - O combate à privação e à violação de direitos ; - A ampliação do acesso a informações sobre a diversidade da nação brasileira e sobre a recriação e afirmação das \textbf{identidades} provocadas pelas relações étnico-raciais ; - A formação continuada dos professores oferecidas nos diferentes etapas e modalidades de ensino. 5.6.3 Ações Educativas de Combate ao \textbf{Racismo} e a \textbf{Discriminações} O princípio encaminha para : - A conexão dos objetivos, de estratégias de ensino e de atividades com a experiência de vida dos alunos e professores, valorizando aprendizagens vinculadas às suas relações com pessoas negras, \textbf{brancas}, mestiças, assim como as vinculadas às relações entre negros, \textbf{indígenas} e \textbf{brancos} no conjunto da sociedade ; - A readequação dos materiais didáticos que promovam o combate ao \textbf{racismo} e as \textbf{discriminações} efetivada pelos profissionais da educação, pelas representações dos negros, dos \textbf{indígenas} e de outras minorias ; - A reflexão e a tomada de decisão acerca das relações étnico-raciais positivas para que professores e alunos possam reconhecer suas responsabilidades enfrentando e superando discordâncias, conflitos, contestações ; - A valorização da oralidade, da corporeidade e da arte, marcas da cultura de \textbf{raiz} africana e \textbf{indígena}, ao lado da escrita e da leitura ; - A educação patrimonial, a partir do patrimônio cultural afro-brasileiro e \textbf{indígena}, visando a preservá -lo e a difundi -lo ; - A ênfase na valorização da participação dos diferentes grupos sociais, étnico-raciais na construção da nação brasileira, aos elos culturais e históricos entre diferentes grupos étnico-raciais ; - A elaboração de projetos políticos pedagógicos que contemplem a diversidade étnico-racial.
Para reafirmar a política curricular para a educação das relações étnico-raciais e quilombolas, faz -se necessário que os entes federados articulem suas ações aos princípios da consciência política e histórica da diversidade ; do fortalecimento de \textbf{identidades} e de direitos pautados em ações educativas de combate ao \textbf{racismo} e a \textbf{discriminações}.
Nessa perspectiva, urge desenvolver ações que busquem superar as práticas discriminatórias étnico-raciais no ambiente escolar, investindo na formação dos profissionais da Educação Básica e na elaboração de materiais didáticos que levem a comunidade escolar a refletir sobre suas práticas pedagógicas na preparação do/a educando/a para o exercício pleno da cidadania, considerando a pluralidade étnico-racial brasileira e atendendo aos dispositivos previstos na Lei de Diretrizes e Bases da Educação Nacional, alterada pela Lei 10.639/2003 ( BRASIL, 2003 ) e Lei 11.645/2008 ( BRASIL, 2008 ) e pelo Plano Nacional de Implementação das Diretrizes Curriculares Nacionais para Educação das Relações Étnico-raciais e para o Ensino de História e Cultura Afrobrasileira e \textbf{Indígena} ( BRASIL, 2009a ).

% matched lemmas: antropológico, branco, contraposição, discriminação, identidade, indígena, negar, popular, racial, racismo, raiz, índio
\end{textsample}
